% gerado por ferramentas/md2tex_enunciado.py a partir de questoes/q118.md (não editar à mão)
A figura ilustra uma forma de medir a densidade relativa de dois líquidos, A e B, imiscíveis, colocando-os em um tubo em forma de U. Atingido o equilíbrio hidrostático, medem-se as alturas, $h_{A1}$ e $h_{B1}$, das colunas dos líquidos, a partir de sua superfície até o nível da linha de separação entre eles. Considere que um volume adicional de líquido A é colocado na coluna esquerda do tubo, conforme a figura. Após o equilíbrio ser restabelecido, medem-se as novas alturas, $h_{A2}$ e $h_{B2}$, da mesma maneira, comparando-as com as alturas anteriores.\par
\figq[0.28\linewidth]{q118_fig1.png}
A comparação entre as novas alturas e as anteriores revela que\par
\begin{altlinha}
\alth[0.325]{A}{$h_{A2} = h_{A1}$ e $h_{B2} = h_{B1}$.}%
\alth[0.325]{B}{$h_{A2} > h_{A1}$ e $h_{B2} = h_{B1}$.}%
\alth[0.325]{C}{$h_{A2} > h_{A1}$ e $h_{B2} > h_{B1}$.}\\[1.5pt]
\alth[0.325]{D}{$h_{A2} < h_{A1}$ e $h_{B2} > h_{B1}$.}%
\alth[0.325]{E}{$h_{A2} = h_{A1}$ e $h_{B2} > h_{B1}$.}
\end{altlinha}
